\scanpage{75}
\unsure{S|Si} Car \deleted{en localisant aux pts génériques}
\[
\dim\bar S_x+\underbrace{\operatorname{codim}(\bar S_x,\bar X_x)}_{\geq m}\leq\dim\bar X_x=n+1
\]
donc par l'hyp. \deleted{de récurrence} $2_{n'}$ pour $n'\leq n+1$
[en particulier $2_{n+1-m}$], on conclut
\[
H_T^{i+1}(\bar S,F)=0\quad\text{si }i>n+1-m.
\]

\uline{Si} $n+1-m\geq2$, \inserted{i.e. $m\leq n-1$,} on a de \unsure{$=$}
\[
H^i(\bar S,F)=H^i(\bar S_0,F_0)\ =\ 0\quad\text{si }i>n+1-m\geq2
\]
d'où $H^i(S,F)=0$. Reste le cas
$m\geq n$, \inserted{donc ($S$ étant $\ne\varnothing$) $m=n$ ou $n+1$} \deleted{Si $m\geq n+1$,} \inserted{\deleted{Si $m=n$, et} Si $m=n+1$,} alors $S$ est \deleted{un pt}
vide ou une \deleted{\unsure{et partie de}} \inserted{partie finie de} $X_0$, et on a
bien $H^i(S,F)=0$ pour $i>0=(n+1)-m$, OK.

Si $m=n$, alors \deleted{le \unsure{seul cas}} \unsure{ch. que|chaque} comp.
\deleted{irréductible \unsure{de}} $S$ est de dim $\leq1$. \deleted{Pour les
\unsure{comp.} irréductibles qui ne sont pas
$=X_0$,} \inserted{Si $S\ne X_0$} on a $\bar S_0$ de dim 0, donc
$H^i(\bar S_0,F)=0$ si $i>0$, à fortiori si $i>n+1-m=1$.

Pour $S=X_0$, on est ramené à prouver
que sur un corps sép. clos \unsure{$k$, ]} son $k\unsure{(t)|[t]|(x)}$ est de \deleted{dim. \unsure{coh|cd}} $\operatorname{cd}_{\ell}\leq1$,
ce qui est connu.
% archive sheet number: 74
